\documentclass[../main.tex]{subfiles}
\begin{document}
%==========================================TIÊU ĐỀ TẬP========================================
\begin{table}[h]
  \begin{adjustwidth}{-1cm}{}
    \begin{tabular}{>{\centering\arraybackslash}p{6cm}|>{\raggedright\arraybackslash}p{12cm}}
      \multirow{5}{6cm}
      % Cột bên trái: ÉPISODE và số tập
      {\\[-40pt]
        \centering
        {\fontsize{20pt}{20pt}\selectfont{\outlinetext[1]{eptype}{white}{\flaregothic SPÉCIAL}}\\[12pt]
          {\fontsize{72pt}{72pt}\selectfont{\outlinetext[1]{epnum}{white}{\burbank XXVIII}}}\\[6pt]       % use for normal/special episodes
          % {\fontsize{36pt}{36pt}\selectfont{\outlinetext[1]{epnum}{white}{\burbank SÁU}}}\\[6pt]     % use for custom episodes
          {\fontsize{20pt}{20pt}\selectfont{\outlinetext[1]{epnum}{white}{\cafeteria (S-28)}}}\color{black}}
        \\[18pt]
        {\fontsize{14pt}{14pt}\selectfont{\outlinetext[1]{epattr}{white}{\flaregothic BIENVENUE, }}{\flaregothic\color{Nodoka} Nodoka}\fontsize{14pt}{14pt}\selectfont{\outlinetext[1]{epattr}{white}{\flaregothic \ !}}}\\
        {\fontsize{18pt}{18pt}\selectfont{\outlinetext[1]{epattr}{white}{\flaregothic POISSON D'AVRIL !}}}\\
        \color{black}
      }
       &
      % Dòng thứ nhất: tên tập tiếng Nhật
      {\fotskip\fontsize{18pt}{18pt}\selectfont \ruby{新}{しん}\ruby{卒}{そつ}で\ruby{入}{はい}った\ruby{事}{じ}\ruby{務}{む}\ruby{所}{しょ}の\ruby{闇}{やみ}がヤバすぎる……} \\[6pt]
      % Dòng thứ hai: tên tập tiếng Anh
       & {\cafeteria\fontsize{18pt}{18pt}\selectfont New Grad Joins Company, Finds Out About Its Dark Side}                                     \\[4pt]
      % Dòng thứ ba: tên tập tiếng Việt
       & {\vnmsans\fontsize{18pt}{18pt}\selectfont Nhân viên học việc}                                                                          \\[8pt]
      % Dòng thứ tư: Nhân vật xuất hiện
       & {\caviar{\fontsize{14pt}{14pt}\selectfont Nhân vật: \emp{kuon2} \emp{game} \emp{achan} \emp{nodoka} \emp{sora1}}}
      % Dòng thứ năm (nếu có): Nhân vật chỉ xuất hiện trong Omake
      %  & {\abv Truyện phụ:} \emp{okayu} 
      \\[8pt]
      % Dòng cuối cùng: Thời gian diễn ra của tập (thường sẽ là ngày phát hành)
       & {\sen\fontsize{16pt}{16pt}\selectfont Mốc thời gian: 01/04/2022}                                                                       \\[18pt]
    \end{tabular}\\[8pt]
  \end{adjustwidth}
\end{table}
%=========================================TRUYỆN CHÍNH========================================
\color{black}

\txtbx[2]{
  {\color{vol} \uline{Tóm tắt:}} Ngày 1 tháng 4, Harusaki Nodoka chính thức gia nhập \hll\ và được Kuon, A-chan cùng Game dẫn đi tham quan. Cả ba liên tục ``hù dọa'' cô bằng những câu chuyện vừa có thật vừa vô lý: văn phòng hay nổ tung, tài năng mang vũ khí hạng nặng, phòng tập có trọng lực gấp ba lần và sóng âm đánh văng người, thậm chí cả chuyện công ty dùng thành viên phi nhân làm lương thực dự phòng. Thay vì sợ hãi bỏ chạy, Nodoka lại tin sái cổ và thể hiện quyết tâm rèn luyện thể chất, học võ, đeo tạ đi làm, khiến ba ``kẻ chủ mưu' vừa buồn cười vừa áy náy. Cuối cùng, Sora xuất hiện với pháo giấy và tiết lộ tất cả chỉ là trò đùa Cá tháng Tư. Nodoka sốc nhưng vẫn tràn đầy nhiệt huyết, và món quà chào mừng thực sự của cô là một núi hồ sơ công việc.
}

Ánh đèn vàng nhạt từ chiếc bóng điện treo lủng lẳng trên trần tỏa ra luồng sáng khiêm tốn, bao phủ lấy căn hầm ngầm quen thuộc. Nơi đây vốn là ``căn cứ bí mật'' mà Kuon cùng nhóm bạn đại học từng dùng để thử nghiệm cuốn sổ Peko quái gở năm nào\footnote{{\flaregothic ÉPISODE 116}, quyển số Chín}. Xung quanh tường, những chiếc kệ sắt chất cao chót vót những chồng sách cũ, các sơ đồ hình tròn cùng công thức bí ẩn được ghim chật kín trên bảng bần, bên cạnh chiếc tivi CRT đời cổ đặt chênh vênh trên chồng thùng carton. Trên góc bảng, tấm lịch treo tường đã được lật sẵn sang tờ ghi rõ chữ ``tháng Tư''.

Trong căn hầm chật hẹp ấy, Sora, A-chan và Kuon đang ngồi quanh chiếc bàn gỗ, rôm rả bàn bạc về một kế hoạch đặc biệt sẽ diễn ra vào ngày mai - ngày mùng 1 tháng Tư - cùng với sự tham gia của Game từ chiếc đồng hồ đeo tay quen thuộc của Kuon.

Sora nhìn tờ lịch trên tường, nhẹ nhàng phá tan sự im lặng: ``Cuối cùng thì ngày này cũng đã tới rồi nhỉ.''

Kuon dựa lưng vào ghế, hai tay khoanh trước ngực, ánh mắt thoáng chút hoài niệm: ``Mọi thứ có vẻ cũng đã được chuẩn bị vô cùng kỹ lưỡng rồi đấy. Tự nhiên làm anh nhớ tới hồi mấy anh em mình từng ngồi lại trong căn hầm này để bàn về kế hoạch quảng bá cho live show đầu tiên\footnote{{\flaregothic ÉPISODE 37}, quyển số Ba của \textit{holo no graffiti}.} của mấy đứa ghê.''

Nàng Thần tượng khẽ mỉm cười, ánh mắt tràn ngập kỷ niệm: ``Hồi đó Miko-chan hoảng loạn tới nỗi sợ hãi luôn cơ mà.''

A-chan đẩy nhẹ gọng kính cận trên chiếc mũi thanh tú, ngón tay chỉ vào bản kế hoạch chi tiết đặt trên bàn: ``Chị với Kuon-san sẽ đảm bảo mọi thứ diễn ra đúng như kịch bản. Còn em thì cứ yên tâm lo phần kết thúc hoành tráng nhé, Sora-chan.''

Sora nở một nụ cười đáp lại. Thế nhưng, khác với nụ cười thiên thần dịu dàng thường ngày, khóe môi cô nàng idol thế hệ đầu tiên khẽ nhếch lên một nét tinh nghịch đầy bí hiểm: ``Ngày mai chắc chắn sẽ thú vị lắm đây\\~{}{}''

Mặt chiếc đồng hồ trên cổ tay Kuon bỗng lóe lên ánh sáng xanh lam, đồ thị sóng âm giật nảy một nhịp như thể đang rùng mình ớn lạnh, trước khi giọng nói trầm khàn của Game cất lên đầy châm chọc: ``Tôi bắt đầu cảm thấy sợ nụ cười của cô rồi đấy, Tokino-san\dots''

Vệt sáng trên mặt kính xoay nhẹ một góc về phía Kuon và A-chan, kèm theo tiếng thở dài điện tử đầy băn khoăn: ``Kuon-user, Yuujin-san, hai người có chắc đây là một ý tưởng hay ho để chào đón cô lính mới sắp gia nhập công ty không đấy?''

A-chan nhếch môi cười khoái chí, giọng điệu đầy tự tin: ``Ngày mai là Cá tháng Tư mà! Mới cả trêu chọc một chút để giải tỏa căng thẳng cho nhân viên mới thì có gì sai đâu chứ?''

Tín hiệu âm thanh từ chiếc đồng hồ khựng lại một nhịp, như thể thực thể A.I kia vừa bị nghẹn lời trước thái độ khác lạ của cô thư ký mẫu mực: ``Tôi ngạc nhiên là những lời đùa cợt này lại được thốt ra từ miệng của một cô gái sống chết bán mình cho tư bản như cô đấy\dots''

Kuon gật gù, cất tiếng giải thích bằng giọng điệu điềm tĩnh nhưng chứa đựng sự trải đời: ``Với cả dù gì, em ấy cũng sẽ phải sớm đối mặt với thực tế công việc `bất bình thường' ở \hll\ thôi. Có thể coi đây là một bài kiểm tra nhập môn\dots\ giúp em ấy làm quen dần. Thà để em ấy nếm trải sự thật ngay từ đầu còn hơn là dỗ dành ngon ngọt rồi sau này lại bị giội một gáo nước lạnh phũ phàng.''

Một tiếng thở dài rè rè như tiếng quạt tản nhiệt vang lên từ loa ngoài của chiếc đồng hồ, thể hiện sự bất lực hoàn toàn trước sự đồng lòng của cả nhóm: ``\dots Cũng phải. Chỉ hy vọng là cô nhân viên mới đó không ngất xỉu vì sốc ngay trong ngày đầu tiên đi làm thôi\dots''

Sora đứng dậy, kéo lại nếp áo ngay ngắn rồi mỉm cười chào tạm biệt mọi người để trở về chuẩn bị cho vai trò quan trọng của mình. Bóng cô khuất dần sau cánh cửa hầm, để lại căn hầm ngầm chìm trong bầu không khí hồi hộp.

A-chan thu dọn lại tập tài liệu trên bàn, còn Kuon thì nhấp một ngụm trà mát lạnh; tất cả đã sẵn sàng cho kịch bản ``Cá tháng Tư'' đầy hỗn loạn sắp sửa giội xuống đầu cô tân binh vào sáng mai.

\begin{center}
  \uline{Ngày hôm sau\\Văn phòng \hll.}
\end{center}

Ánh nắng ấm áp của ngày mùng 1 tháng Tư chiếu qua ô cửa kính, làm bừng sáng cả gian phòng làm việc của trụ sở \hll.

Bên trong văn phòng, đối diện với Kuon và A-chan lúc này là một cô gái trẻ trong bộ trang phục công sở phẳng phiu, chỉn chu. Cô đứng thẳng lưng, hai tay chắp trước bụng, nhắm mắt hít một hơi thật sâu rồi chậm rãi thở ra để lấy lại dũng khí. Sự hồi hộp của một tân binh trong ngày đầu đi làm hiện rõ trên từng cử chỉ.

Kuon mỉm cười nhẹ nhàng, lên tiếng phá tan sự ngột ngạt: ``Được rồi, bây giờ thì em hãy giới thiệu bản thân đi.''

Nữ nhân viên đứng nghiêm cẩn, ánh mắt tràn đầy quyết tâm giương về phía hai người đồng nghiệp cấp trên tương lai của mình. Cô cất giọng trong trẻo, dõng dạc: ``Chào buổi sáng, hai anh chị! Em là Harusaki Nodoka, và em sẽ bắt đầu làm việc với mọi người từ hôm nay!''

\chrint

Đó chính là Harusaki Nodoka (\ruby{春}{はる}\ruby{先}{さき}のどか, \ruby{Ha-ru}{Xuân}-\ruby{sa-ki}{Tiên}\ Nô-đô-ca), nữ nhân viên tập sự mới toanh vừa chính thức gia nhập đại gia đình \hll. Là hậu bối trực tiếp dưới trướng của A-chan và Kuon, cô gái trẻ này mang trong mình nguồn năng lượng dồi dào cùng khát khao cháy bỏng được cống hiến cho công việc văn phòng và hỗ trợ các nữ thần tượng.

\model{25}{l}{6.5cm}{nodoka}

Mục tiêu ngắn hạn hàng đầu của Nodoka chính là tiếp bước người tiền bối dẫn dắt mình, nỗ lực học hỏi từng chút một để sớm trở thành một nhân viên mẫn cán, hữu ích cho công ty. Thế nhưng, ẩn sau dáng vẻ đầy tinh thần trách nhiệm và tác phong luôn cố gắng tỏ ra chuyên nghiệp ấy, cô nàng thực chất lại là một kẻ có đôi chút đãng trí và lơ đãng. Đặc biệt, Nodoka còn sở hữu một niềm đam mê mãnh liệt, vô cùng cuồng nhiệt với\dots\ loài ếch.

Xét về tính cách, Nodoka là hiện thân của một tân binh điển hình: ngoan ngoãn, lễ phép, chân thành và cực kỳ nghiêm túc trong mọi việc được giao. Mỗi lần chào hỏi, cô luôn giữ thói quen cúi gập người đúng 90 độ chuẩn chỉ, ăn nói bằng kính ngữ mực thước và luôn lắng nghe mọi lời chỉ bảo với đôi mắt sáng ngời. Thế nhưng, chính sự nghiêm túc đến mức ngây thơ cùng nét tính cách hơi ``PON'' (vụng về) bẩm sinh lại khiến cô nàng thường xuyên rơi vào những tình huống dở khóc dở cười. Mỗi khi gặp sự cố bất ngờ hay bị các tiền bối trêu chọc, Nodoka rất dễ luống cuống, mặt đỏ bừng lúng túng, tạo nên một nét đáng yêu vô cùng trong trẻo khiến ai nấy đều muốn chở che.

Về ngoại hình, Nodoka toát lên vẻ đẹp thanh tú, dịu dàng đúng như cái tên mang ý nghĩa ``tiết trời xuân thanh bình'' của mình. Cô sở hữu mái tóc ngắn cắt kiểu bob màu nâu hạt dẻ bồng bềnh, điểm xuyết một chiếc kẹp tóc nhỏ màu xanh da trời kẹp gọn ghẽ bên mái phải. Đôi mắt to tròn mang sắc nâu hạt dẻ phảng phất ánh hồng long lanh như biết cười, kết hợp cùng làn da trắng mịn và nụ cười tươi tắn luôn thường trực trên môi, mang lại cho người đối diện cảm giác vô cùng ấm áp, dễ chịu.

Bộ trang phục thường ngày của cô mang đậm phong cách của một nữ nhân viên công sở trẻ trung và chỉn chu. Bên trong là chiếc áo sơ mi trắng có cổ bẻ gọn gàng, khoác bên ngoài chiếc áo len cardigan mỏng cài cúc màu xanh bơ nhạt mát mắt với phần tay áo xắn nhẹ năng động. Phía dưới, cô diện chiếc chân váy bút chì màu đen ôm dáng thanh lịch, đi cùng đôi quần tất mỏng màu đen tuyền và đôi giày cao gót đế thấp màu đen bóng trang nhã. Nổi bật nhất trước ngực áo của cô là chiếc dây đeo thẻ nhân viên màu hồng in họa tiết đặc trưng của \hll, giữ chiếc thẻ nhân viên viền hồng có đính kèm một chiếc huy hiệu hình chú ếch xanh ngộ nghĩnh - nét chấm phá đáng yêu khẳng định sở thích độc nhất vô nhị của cô nàng học việc.

\chrint

Cô nàng chỉnh lại chiếc thẻ nhân viên trước ngực, tiếp tục phần chào hỏi bằng một năng lượng tràn đầy: ``Em sẽ cố gắng hết sức ạ! Mong hai anh chị hãy chỉ giáo cho em!'' Nodoka kết thúc phần giới thiệu bằng một cái gập người 90 độ chuẩn chỉnh.

A-chan và Kuon nhìn cô nàng tân binh tóc nâu tràn đầy sức sống ấy, rồi quay sang nhìn nhau khẽ gật đầu đồng tán thưởng.

Nàng Đeo kính mỉm cười thân thiện: ``Ồ, vậy em là Nodoka-san à. Chào buổi sáng em nhé.''

Kuon cũng gật đầu chào đón: ``Ừm. Chào mừng em tới với \hll.''

Nodoka ngẩng đầu lên, hai mắt sáng rực niềm vui: ``Dạ, A-chan-san và Kuon-san! Em rất vui vì được làm việc với hai người!''

Nhìn nét mặt rạng rỡ, ngây thơ và tươi tắn của cô nhân viên mới, A-chan không thể kìm lòng được trước sự đáng yêu của người đàn em. Cô áp hai tay lên má, thốt lên đầy phấn khích: ``Ối chà, cuối cùng thì chị cũng có một người đàn em dễ thương rồi!''

Nodoka lúng túng, hai má đỏ hây hây, xua tay bối rối: ``D-Dạ\dots\ Em cũng không dám nhận đâu ạ\dots''

Kuon thấy A-chan sắp sửa ``làm quá'' liền lên tiếng cắt ngang: ``Bình tĩnh lại nào, A-san, chị đang làm em ấy sợ đấy.'' Rồi, cậu quay sang nhìn cô nàng tân binh đang ngượng ngùng kia, ``trông em cũng có vẻ rất lịch sự và chuyên nghiệp đấy, Nodoka-chan.''

Nodoka cúi đầu cảm kích: ``E-Em cảm ơn anh ạ, Kuon-san!''

Cậu Tổng đưa tay xem giờ, thong thả nói tiếp: ``Dù gì thì, em cũng vừa mới chân ướt chân ráo tới đây, có lẽ ta cũng nên bắt đầu tham quan một vòng nơi này thôi nhỉ?''

A-chan vỗ tay hào hứng hưởng ứng: ``Ý hay đấy! Chúng ta bắt đầu ngay nhé. Chị sẽ dẫn em đi dạo một vòng quanh văn phòng này trước tiên.''

Nodoka cúi đầu, mỉm cười ngoan ngoãn: ``Dạ vâng! Em làm phiền hai anh chị rồi ạ.''

Đúng lúc cả ba vừa quay người bước đi, mặt chiếc đồng hồ đeo tay thông minh trên cổ tay trái của Kuon bỗng lóe lên một tia sáng xanh nhẹ.

Từ chiếc loa nhỏ xíu phát ra một giọng nam trầm ấm, tự nhiên và vô cùng rõ ràng như thể có một người thật đang đứng sát bên cạnh trò chuyện: ``Này, Yuujin-san, Kuon-user. Trước khi dẫn lính mới đi tham quan, hai người có nhớ dặn cô ấy cẩn thận với cái máy bán hàng tự động hay nuốt người ở cuối hành lang hay mấy cái góc tường hay tự động biến đổi cấu trúc chưa đấy?''

Nodoka đang bước đi bỗng khựng lại như trời trồng. Cô giật mình nảy người, hai mắt trợn tròn, kinh hãi nhìn chằm chằm vào cổ tay trái của Kuon rồi lắp bắp không thành tiếng: ``Ể!? V-Vừa rồi… có tiếng người phát ra từ… từ trong chiếc đồng hồ của Kuon-san kìa!?''

Cô hoảng hốt lùi phắt lại hai bước, hai tay che miệng, chỉ tay run run vào cổ tay Kuon cứ như thể vừa thấy một vật thể lạ ngoài hành tinh: ``Ch-Chẳng lẽ chiếc đồng hồ đó bị yểm bùa hay chứa linh hồn bên trong hả anh!?''

A-chan bật cười khúc khích trước vẻ mặt hốt hoảng tột độ của cô tân binh, dịu dàng bước tới vỗ vai trấn an: ``À, em đừng sợ quá, Nodoka-san. Đó chỉ là\dots một `cư dân đặc biệt' gắn liền với bọn chị thôi.''

Kuon tỏ vẻ thản nhiên, tiện tay khẽ chạm vào mặt đồng hồ trên cổ tay: ``Đúng thế. Cứ coi như đây là một tính năng hỗ trợ ảoảo\dots\ kiểu trợ lý thông minh đời mới đi. Chỉ có điều trợ lý này hơi nhiều chuyện và thích xen vào chuyện người khác một chút thôi.''

Giọng nói của Game từ trong chiếc đồng hồ lập tức vang lên phản bác, mang theo tiếng cười đầy vẻ bất đắc dĩ nhưng vô cùng sống động: ``Này này, ai là 'trợ lý ảo' nhiều chuyện hả, Kuon-user? Tôi nghe thấy hết đấy nhé! Dù sao thì, chào mừng cô đến với \hll, Harusaki-san. Chúc cô giữ được cái biểu cảm hoảng hốt tràn đầy sức sống đó suốt cả ngày hôm nay nha.''

Nodoka đứng chết trân tại chỗ, mồ hôi hột túa ra như mưa trên trán. Trong đầu cô lúc này hàng vạn câu hỏi chấm đang bay lượn kèm theo một sự hoang mang tột độ: ``\textit{Đồng hồ đeo tay biết nói chuyện dõng dạc, lại còn biết trêu chọc người khác\dots\ Rốt cuộc đây là công ty giải trí thần tượng hay là viện nghiên cứu dị dạng khoa học viễn tưởng vậy giời!?}''

Ngay khi hai cô gái vừa đi trước một đoạn, mặt chiếc đồng hồ đeo tay thông minh trên cổ tay trái của Kuon chợt lóe lên ánh sáng xanh nhẹ.

Từ loa nhỏ của chiếc đồng hồ, giọng nói của Game vang lên: ``\hl{Thái độ ngây thơ chuẩn chỉnh của một tân binh\dots{} Nhìn cô ấy hào hứng và tràn đầy nhiệt huyết thế này, tôi lại thấy hơi cắn rứt lương tâm cho cái bẫy Cá tháng Tư mà chúng ta sắp giội xuống đầu cô ấy rồi đấy, Kuon-user.}''

Kuon khẽ đưa cổ tay lên gần miệng, hạ thấp giọng thì thầm: ``\hl{Tôi cũng thấy vậy đấy, nhưng biết sao được, kế hoạch là vậy mà.}''

A-chan hơi ngoảnh đầu lại phía sau, khẽ nhếch môi trêu chọc chiếc đồng hồ trên tay Kuon: ``Bớt đóng vai người tốt đi Game-san. Tối qua ai là người gợi ý thêm mấy tình huống giật gân để `thử thách lòng tin' của em ấy vậy?''

Giọng Game từ trong chiếc đồng hồ phát ra một tiếng thở dài nhè nhẹ, tràn đầy vẻ bất lực nhưng cũng không giấu nổi sự buồn cười: ``Tôi chỉ đóng góp ý kiến mang tính chuyên môn thôi, Yuujin-san\dots\ Dù sao thì, nhìn biểu cảm của cô lính mới này kìa. Hy vọng kịch bản của hai người không làm cô ấy sợ đến mức nộp đơn xin nghỉ việc ngay trong ngày đầu tiên đấy nhé.''

Kuon mỉm cười bí hiểm, chỉnh lại gấu tay áo che đi chiếc đồng hồ rồi sải bước đuổi theo phía sau: ``Yên tâm đi. Mọi thứ chỉ mới bắt đầu thôi\dots''

\ct

Cả ba dẫn Nodoka đi dọc theo hành lang rộng rãi, tiến về phía một khu vực cách đó không xa. Đập vào mắt Nodoka là một không gian vô cùng êm ái với chiếc ghế sô-pha dài bọc da cao cấp, chiếc TV màn hình lớn dựng ngay trung tâm và một chiếc tủ kính chất đầy các loại bánh kẹo, đồ ăn vặt hấp dẫn.

A-chan dừng lại, dõng dạc giới thiệu như một hướng dẫn viên chuyên nghiệp: ``Đây là khu vực nghỉ ngơi dành cho các tài năng. Nơi này được thiết kế để họ thư giãn và nạp lại năng lượng sau những giờ làm việc căng thẳng. Nhưng lưu ý nhé\dots\ nơi này rất hay nổ tung, vì vậy em hãy luôn giữ cảnh giác cao độ.''

Nodoka ngoan ngoãn gật đầu ghi nhớ: ``Dạ, em hiểu rồi ạ\dots'', nhưng rồi đột nhiên cô khựng lại giữa chừng. Đôi mắt nâu tròn xòe trợn ngược, cô quay phắt lại nhìn A-chan Rồi cô đột nhiên khựng lại giữa chừng, đôi mắt nâu tròn xòe trợn ngược: ``Khoan đã! A-chan-san\dots\ Nổ tung ấy ạ!?''

Mặt đồng hồ trên tay Kuon lại phát ra ánh sáng xanh nhè nhẹ. Giọng nói trầm ấm của Game từ loa nhỏ vang lên, vô cùng điềm tĩnh tới cô nàng tân binh vào hoảng loạn: ``Cô không nghe nhầm đâu, Harusaki-san. Là `nổ tung' theo đúng nghĩa đen, kèm theo sóng xung kích, khói bụi mù mịt và gạch đá bay tứ tung đấy.''

A-chan khoanh tay trước ngực, gật đầu xác nhận với ánh mắt cực kỳ nghiêm túc: ``Đúng vậy. Mọi người ở đây khi làm việc thường tiêu hao rất nhiều năng lượng. Năng lượng đó tích tụ dần theo thời gian và đôi lúc sẽ bị mất kiểm soát rồi bùng nổ.''

Nodoka nuốt nước bọt ừng ực, hai tay ôm lấy ngực, ánh mắt vừa kinh hãi vừa bàng hoàng: ``E-Em chưa bao giờ nghĩ rằng các tiền bối lại sở hữu nguồn năng lượng dồi dào và bùng nổ đến mức thế\dots''

Mặt đồng hồ lại lóe sáng, vị Trợ lý ảo chêm vào một câu tỉnh rụi: ``Ừ, dồi dào tới mức đôi khi chỉ vì hai cô gái cãi nhau xem ai lỡ ăn mất cái bánh pudding trong tủ lạnh là cả cái văn phòng này liền biến thành tro bụi.''

Kuon thở dài một hơi đầy ngao nán, tay giơ lên xoa xoa thái dương như thể đang sống lại những ký ức đau thương: ``Nói tới mới nhớ, cái phòng này đã tiêu tốn của anh cả núi tiền ngân sách lẫn công sức chỉ để dọn dẹp và xây dựng lại đấy. Mãi gần đây anh chịu hết nổi, mới phải cắn răng lắp đặt hệ thống tự động phục hồi cấu trúc công trường thì mới kịp tiến độ giao việc cho công ty.''

A-chan chỉnh lại gọng kính, khẽ thở dài đồng cảm rồi quay sang nhìn Nodoka: ``Chị cũng từng bị cuốn vào các vụ nổ ở đây rồi, không chỉ một mà tới hai lần lận.''

``Hả!? H-Hai lần lận ạ!?'' nàng tân binh thốt lên, hai đầu gối bắt đầu run rẩy.

Kuon nhếch môi cười nhạt, xua xua tay: ``Chị ấy thế là còn ít đấy. Anh đây còn bị dính trực diện ít nhất bốn, năm quả nổ tung ngay trước mặt cơ. May mà cái xác này còn dẻo dai giữ được mạng, chứ không là lên bàn thờ ngắm gà khỏa thân từ đời tám hoánh nào rồi.''

Nàng Đeo kính gật gù, kết luận một câu đầy sức nặng: ``Thế nên nếu một nhân viên không có một cơ thể dẻo dai chịu đựng được áp lực từ những vụ nổ thì khó mà sống sót, ý chị là khó mà làm việc lâu dài ở đây được, Nodoka-san ạ.''

Cậu Tổng chêm thêm một ngón tay lên: ``Hoặc là không có đôi chân đủ nhanh để đào tẩu và một cái óc tính toán quỹ đạo mảnh vỡ thật chuẩn thì cũng chả đỡ nổi đâu.''

Nodoka hoàn toàn sốc ngang. Cô đứng chết trân tại chỗ, linh hồn như muốn bay ra khỏi khuôn miệng đang há hốc. Trong đầu cô tân binh lúc này là viễn cảnh bản thân phải vừa né bom đạn vừa bế tài liệu chạy thục mạng mỗi ngày.

Thấy vẻ mặt như sắp ngất đi của đàn em, A-chan liền nhẹ nhàng bước tới, vỗ nhẹ lên vai cô mỉm cười an ủi: ``Đừng lo lắng quá nhé, Nodoka-san. Rồi em sẽ quen với tần suất nổ ở đây thôi.''

Kuon cũng gật đầu đồng tình: ``Mọi chuyện rồi đâu sẽ lại vào đấy mà. Cùng lắm là đi bó bột vài tuần.''

Thế nhưng, trái ngược với suy đoán của hai người rằng cô tân binh sẽ hoảng sợ mà nộp đơn xin nghỉ, Nodoka đột ngột siết chặt hai nắm đấm trước ngực. Đôi mắt nâu hạt dẻ của cô bùng cháy lên ngọn lửa quyết tâm ngút ngàn, cô gập người 90 độ dũng cảm hô lớn: ``D-Dạ, em hiểu rồi ạ! Em nhất định sẽ rèn luyện thể chất thật tốt, tập thể dục mỗi ngày để sớm có một cơ thể siêu dẻo dai, đủ sức gánh vác và giúp ích cho mọi người ạ!''

Nói rồi, cô hào hứng tiến lại gần chiếc tủ đồ ăn vặt để quan sát kỹ hơn ``hiện trường nguy hiểm''.

Nhìn tấm lưng tràn đầy sinh lực và quyết tâm hừng hực của cô tân binh, ba ``kẻ chủ mưu'' vội chụm đầu lại với nhau, hạ thấp giọng thì thầm bàn bạc.

A-chan lấy tay che miệng, thì thào với vẻ hoảng hốt: ``\hl{Này Kuon-san\dots{} Con bé tin sái cổ luôn kìa! Chị thấy hình như bọn mình chém gió hơi quá đà rồi đấy?}''

Kuon nhún vai đáp lại: ``\hl{Chém gì đâu chị. Mình đang nói sự thật đấy chứ. Chẳng qua là chúng ta đang hơi phóng đại một chút thôi.}''

Mặt đồng hồ của Kuon khẽ rung lên, giọng Game thì thầm đầy e ngại: ``\hl{Tôi bắt đầu cảm thấy sự thật thà và lòng quyết tâm của cô ấy là một mối đe dọa ngược lại cho chúng ta rồi đấy\dots{} Nhỡ cô ấy đi đăng ký lớp huấn luyện đặc công thật thì sao?}''

Kuon gãi gãi đầu, méo mặt thì thầm trả lời hai người: ``\hl{Bây giờ mà thú nhận là bọn mình đùa Cá tháng Tư thì còn ngượng hơn nữa\dots{} Thôi cứ để em ấy tự tập thể dục đi, dù sao thì tập thể dục cũng tốt cho sức khỏe mà.}''

\ct

Rời khỏi khu vực phòng nghỉ, A-chan và Kuon tiếp tục dẫn Nodoka bước đi dọc theo dãy hành lang cách âm yên tĩnh. Cả nhóm dừng lại trước một cánh cửa gỗ dày dặn, phía trên có gắn chiếc đèn led màu đỏ đang hiển thị dòng chữ ``Đang thu âm''.

Nàng Đeo kính nghiêm mặt, xoay người lại nghiêm cẩn dặn dò đàn em: ``Đây là phòng thu âm của công ty. Hãy cực kỳ cẩn thận mỗi khi bước chân vào đây nhé.''

Nàng Tân binh chớp mắt, cố gắng dùng tư duy logic của một nhân viên công sở bình thường để suy đoán: ``Dạ? Có phải vì hệ thống micro bên trong quá nhạy, nên nếu không cẩn thận thì tiếng bước chân hay tiếng thở của mình sẽ tự động bị thu âm lại ạ?''

Cậu Tổng lắc đầu, xua xua tay với một nụ cười đầy ái ngại: ``Ờm, không đâu. Nếu mà chỉ đơn giản như em nói thì đã chẳng có điều gì đáng lo cả. Đằng này thì\dots''

A-chan đẩy nhẹ gọng kính, ánh mắt trở nên vô cùng bí hiểm và nghiêm trọng: ``Mọi thứ còn phức tạp và mang tính 'sinh tồn' hơn em nghĩ nhiều. Các tài năng có trang bị vũ khí có thể bất ngờ bộc phát tấn công em một cách hoàn toàn không báo trước đấy.''

Chiếc đồng hồ trên cổ tay Kuon lập tức phát ra ánh sáng xanh nhè nhẹ. Giọng Game cất lên trầm ngâm, liệt kê từng nhân vật với thái độ tỉnh rụi: ``Để tôi lấy ví dụ cho cô dễ hình dung nhé, Harusaki-san\dots\ Xem nào, Nakiri-san thì lúc nào cũng lăm lăm hai thanh katana sắc lẹm giắt sau lưng này, rồi Shirogane-san với cái chùy thép khổng lồ đập đâu phá đấy này\dots''

Kuon chêm thêm một câu bằng giọng điệu không thể thản nhiên hơn: ``Chưa kể tới một vài thành viên khác lúc nào cũng thủ sẵn súng ngắn, súng hạng nặng với cả bom nổ chậm trong người nữa cơ.''

Nodoka nghe xong mà muốn rụng rời tay chân. Cô trợn tròn mắt, hai tay ôm lấy đầu, gào lên trong sự hoảng hốt tột độ: ``\textbf{T-Tại sao học lại mang theo vũ khí hạng nặng vào công ty mà không ai cấm cơ chứ!?} Và quan trọng nhất, họ dùng mấy thứ đồ sát thương đó trong phòng thu âm để làm cái gì vậy ạ!?''

A-chan nhún vai đầy bất lực, làm vẻ mặt chịu trận: ``Chị cũng chẳng biết mục đích thực sự của họ đâu. Nhưng tóm lại là nơi này tiềm ẩn đủ mọi loại nguy hiểm rình rập, vì vậy hãy chắc chắn rằng em đã chuẩn bị tâm lý và trang bị phòng hộ thật kỹ lưỡng trước khi đặt chân vào.''

Nodoka đứng ngơ ngác mất vài giây khi vừa bồi thêm một cú sốc nữa từ đàn anh đàn chị. Nhưng rồi, thay vì hoảng sợ rút lui, ngọn lửa quyết tâm dường như lại bùng lên mãnh liệt hơn trong mắt cô tân binh.

Cô lấy ra chiếc sổ tay nho nhỏ, vừa cẩn thận ghi chép vừa gật đầu lia lịa: ``Em hiểu rõ rồi ạ! Để đảm bảo an toàn tuyệt đối cho bản thân và công việc, ngay cuối tuần này em sẽ ra cửa hàng đi săn một bộ giáp sắt toàn thân hoàn chỉnh để mặc đi làm ạ!''

Nói rồi, nàng tân binh hào hứng tiến lại gần sát cánh cửa phòng thu, nghiêm túc ngó qua ô kính nhỏ để ``thám sát tình hình chiến sự'' bên trong.

Nghía thấy cô tân binh đang mải mê quan sát, ba ``kẻ chủ mưu'' lại lập tức chụm đầu lại với nhau, hạ thấp giọng thì thầm.

Mặt đồng hồ của Kuon rung nhẹ, giọng Game thì thào đầy kinh ngạc lẫn buồn cười: ``\hl{Bộ giáp sắt toàn thân hoàn chỉnh\dots? Cô ấy tính hóa thân thành kị sĩ trung cổ mặc giáp đi làm công sở thật đấy à?}''

A-chan che miệng, mồ hôi hột chảy dài trên trán: ``\hl{Kuon-san ơi\dots{} Chị thấy con bé vừa vẽ xong cái phác thảo bộ giáp chiến đấu vào sổ tay thật kìa! Bọn mình có nên dừng lại không chứ, chị thấy tội lỗi quá\dots}''

Kuon đỡ trán, thở dài một hơi dở khóc dở cười: ``\hl{Ai mà ngờ cái tính thật thà với độ tập trung của em ấy lại cao đến mức này cơ chứ\dots{} Kịch bản này bắt đầu đi lệch hướng và vượt quá tầm kiểm soát của chúng ta rồi đấy\dots}''

\ct

Sau chuyến tham quan phòng thu âm đầy căng thẳng, đồng hồ trên tường cũng vừa điểm 12 giờ trưa. Cả ba dẫn Nodoka trở lại khu vực văn phòng chính để nghỉ ngơi.

A-chan lấy ra vài hộp bento nóng hổi đặt lên bàn, nhẹ nhàng mỉm cười với cô tân binh: ``Chúng ta nghỉ trưa thôi nào. Nếu em thích thì em có thể ăn hộp cơm tonkatsu (thịt heo chiên xù) này nhé.''

Kuon dựa lưng vào thành ghế, mỉm cười hào phóng: ``Anh tự tay vào bếp làm hết đấy. Em cứ ăn thoải mái, không phải ngại đâu.''

Nodoka mở nắp hộp bento ra, mùi thơm lừng của thịt chiên xù quyện với nước sốt đậm đà tỏa ra ngào ngạt khiến bụng cô reo lên một tiếng nho nhỏ. Mắt cô sáng rực: ``Ô, bữa trưa trông ngon và hấp dẫn quá\dots\ Cảm ơn Kuon-san và A-chan-san rất nhiều vì đã đãi em ạ!''

Co nàng dùng đũa gắp một miếng tonkatsu vàng ươm cho vào miệng. Ngay lập tức, lớp vỏ xù giòn rụm vỡ ra, phần thịt bên trong mềm mại, mọng nước ngọt ngào bùng nổ trên đầu lưỡi.

Cô nhai ngấu nghiến, hai má phồng phông lên như sóc ôm hạt, vừa ăn vừa thốt lên trong sung sướng: ``Hộc\dots\ Ngon quá đi mất ạ! Lớp vỏ giòn tan mà thịt bên trong lại mọng nước thơm phức luôn! Kuon-san nấu ăn đỉnh thực sự!''

A-chan nhìn Nodoka ăn một cách ngon lành, chợt ho nhẹ một tiếng rồi đổi sang tông giọng nghiêm túc, sâu lắng: ``Nhân tiện thì, Nodoka-san này\dots\ Theo em thì ưu tiên hàng đầu của công ty chúng ta là gì?''

Nodoka dừng đũa, ngẫm nghĩ một hồi lâu, lông mày khẽ nhíu lại tỏ vẻ rất tâm huyết: ``Ê-tồ\dots\ Chắc là việc giám sát chặt chẽ và điều phối hoạt động của các tài năng ạ?''

Kuon gật gù, điềm nhiên lên tiếng: ``Không sai đâu, đó cũng là một trong những điều vô cùng quan trọng để \hll\ có thể vận hành một cách bài bản. Nhưng đó vẫn chưa phải là ưu tiên hàng đầu tối thượng đâu em.''

Nodoka ngơ ngác, chớp chớp mắt hỏi: ``Vậy\dots\ nó thực sự là cái gì ạ?''

A-chan chống tay lên bàn, ghen tị nhìn hộp bento rồi nghiêm trọng hóa vấn đề: ``Mặc dù điều em ấy nói chắc chắn rất quan trọng, nhưng ưu tiên hàng đầu của chúng ta là phải luôn chuẩn bị sẵn sàng cho các tình huống khủng hoảng có thể xảy tới bất kỳ lúc nào.''

``K-Khủng hoảng ấy ạ?'' Nodoka nuốt ậm ừ miếng thịt, bắt đầu cảm thấy có gì đó sai sai.

A-chan đưa một ngón tay lên lấy ví dụ: ``Chẳng hạn nhé, giả sử như một ngày nào đó có ai đó quên đặt bữa trưa cho toàn bộ các tài năng\dots''

Kuon chêm vào bằng giọng điệu u ám: ``\dots và tất cả mọi người đều rơi vào cảnh cháy túi, hoặc là quá bận rộn đến mức chẳng một ai có thể ra ngoài mua đồ ăn.''

Nodoka hoảng hốt, hai tay ôm lấy má: ``T-Thế thì nguy hiểm quá! Vậy\dots\ chuyện gì sẽ xảy ra tiếp theo ạ!?''

A-chan nhếch môi nở một nụ cười đầy bí hiểm, thản nhiên thả một cú bọc lót chấn động: ``Đó chính là lúc mà \hll\ kích hoạt chính sách đặc biệt: có một vài tài năng được dành riêng để làm\dots\ lương thực dự phòng.''

Nodoka đứng hình mất năm giây, đôi đũa trên tay run rẩy: ``\dots Dạ? Ý của chị là\dots''

Mặt đồng hồ trên tay Kuon lại phát ra ánh sáng xanh nhè nhẹ. Giọng Game cất lên chậm rãi, liệt kê từng cái tên bằng một thái độ vô cùng ``chuyên nghiệp'': ``Có lẽ cô chưa để ý, nhưng trong cái văn phòng này có rất nhiều thành viên thuộc chủng tộc phi nhân đấy. Như là Tsunomaki-san - một con cừu béo mầm này, Usada-san - một con thỏ tăng động này, hay Takanashi-san\footnote{Sẽ được giới thiệu ở quyển Mười Sáu.} - một con gà chiên giòn này\dots''

Kuon nhấp một ngụm trà mát lạnh, bình thản chêm vào: ``Điều đó cho phép ban quản lý chúng ta tiếp cận với một nguồn cung cấp nguyên liệu tươi ngon và phong phú cực kỳ.''

A-chan gật đầu tán thưởng:

``Và thế là, chúng ta vẫn có thể chuẩn bị một bữa ăn khá thịnh soạn để cứu vớt cả công ty khỏi nạn đói.''

Nodoka cúi đầu nhìn xuống miếng tonkatsu dở dang trên tay mình. Mồ hôi hột túa ra như mưa, mặt cô cắt không còn một giọt máu, giọng nói run rẩy đến tội nghiệp: ``V-Vậy\dots\ đừng nói với em rằng miếng thịt chiên mà em đang ăn dở đây là\dots!''

Kuon khẽ thở dài, khoanh tay trước ngực, buông một câu triết lý sống đầy u tối: ``Không cần phải buồn hay dằn xót làm gì đâu em. Đây chỉ là cách mà xã hội vận hành nói chung thôi. Chúng ta đang sống trong một thế giới nơi cá lớn nuốt cá bé, mạnh được yếu thua mà, đúng chứ?''

Nodoka hoảng loạn nội tâm tột độ. Hai mắt cô rơm rớm nước mắt, đầu óc quay mòng mòng: ``\textit{T-Thế giới của người lớn thật là điên rồ và tàn nhẫn\dots\ Và những ràng buộc khắc nghiệt trong một công ty quản lý tài năng cũng vậy\dots!}''

Cô nhìn chằm chằm vào đĩa thức ăn. Nhưng rồi, thay vì nôn ra hay hoảng sợ bỏ chạy, ngọn lửa sinh tồn tột cùng lại bùng lên trong mắt cô tân binh.

Nodoka cắn chặt răng, nuốt nước mắt vào trong, siết chặt đôi đũa với một sự quyết tâm bi ai: ``Em\dots\ Em hiểu rồi! Em không thể lãng phí sự hy sinh cao cả của các tiền bối được! Em sẽ ăn hết sạch không để thừa một hạt cơm nào để đảm bảo sự sống còn của mình trong cái xã hội khắc nghiệt này!''

Nói rồi, cô vừa sụt sịt khóc vừa cắm cúi nhai miếng tonkatsu một cách điên cuồng và bi thảm.

Trong lúc Nodoka đang mải mê ``chiến đấu vì sự sống còn'' trong nước mắt, ba ``kẻ chủ mưu'' vội vã chụm đầu lại với nhau ở góc phòng, hạ thấp giọng thì thầm trong sự bàng hoàng.

A-chan lấy tay che miệng, mồ hôi chảy ròng ròng, thì thào đầy hoảng hốt: ``\hl{Trời ơi Kuon-san! Cái kịch bản `ăn thịt thành viên phi nhân' mà hai chúng mình tự nghĩ ra để hù con bé\dots{} nó tin sái cổ luôn kìa! Em ấy vừa ăn vừa khóc thương cho Watame-san kìa!}''

Kuon gãi đầu, mặt méo xệch thì thầm lại: ``\hl{Ai mà biết được em ấy lại thật thà tới mức ngây ngô này chứ! May mà cái món tonkatsu này là thịt lợn xịn 100\% em mua ở siêu thị, chứ công ty mình đã có thành viên nào là lợn đâu!}''

Mặt đồng hồ của Kuon rung nhẹ, giọng Game thì thầm phát ra, mang theo sự đánh giá vô cùng ngao ngán lẫn cạn lời: ``\hl{Tôi thực sự bó tay toàn tập với độ tin người đến mức mù quáng của cô tân binh này rồi đấy, Yuujin-san, Kuon-user\dots{} Xét về độ ngây thơ đần đệch, tôi đánh giá cô ấy còn vượt mặt cả cô em út Naomi nhà cậu rồi đấy, Kuon-user. Cô bé tuy nhiều lúc ngơ ngơ ngác ngác thật, nhưng còn lâu cô ấy mới tin cái chuyện công ty đem nhân viên ra làm thịt để ăn trưa thế này!}''

Kuon thở dài đỡ trán: ``\hl{Thôi xong rồi\dots{} Tình hình này mà không dừng lại nhanh thì có khi em ấy đi báo cảnh sát bắt bọn mình vì tội ngược đãi động vật thật mất\dots}''

\ct

Sau khi giải quyết xong bữa trưa ``đầy nước mắt và sự hy sinh'', A-chan và Kuon tiếp tục dẫn Nodoka đi tới cuối dãy hành lang tầng hai. Cả nhóm dừng lại trước một cánh cửa đôi bằng kính cường lực cực kỳ kiên cố, bên trên có tấm bảng inox ghi rõ dòng chữ: ``Phòng tập luyện vũ đạo''.

A-chan đẩy nhẹ cánh cửa, hé mở một không gian rộng lớn với hệ thống gương soi kín tường và sàn gỗ bóng loáng. Cô quay lại, dõng dạc giới thiệu: ``Tiếp theo chính là phòng tập luyện vũ đạo. Đây là nơi các tài năng mài dũa kỹ năng trình diễn và chuẩn bị cho các đêm live show hoành tráng.''

Nodoka ngó đầu vào bên trong, hai mắt sáng rỡ: ``Ồ! Phòng tập rộng rãi và hiện đại quá ạ! Nhìn chuyên nghiệp thật đấy!''

Kuon bước tới đứng bên cạnh, lắc đầu nhẹ thở dài, tạo một khoảng lặng đầy nguy hiểm: ``Nhìn bề ngoài có vẻ bình thường thế thôi em ạ\dots\ Chứ thực chất, đây lại là 'vùng nguy hiểm cấp độ S' của toàn bộ cái trụ sở này đấy.''

Nodoka nụ cười cứng đờ trên môi, chớp chớp mắt: ``D-Dạ!? Vùng nguy hiểm cấp S ạ!?''

A-chan chỉnh lại gọng kính cận, cất giọng trầm thấp đầy đe dọa: ``Đúng vậy. Em có biết làm sao để các idol của chúng ta có thể vừa hát vừa nhảy sung sức suốt 3 tiếng đồng hồ trên sân khấu mà không hụt hơi không? Đó là nhờ hệ thống giả lập môi trường khắc nghiệt trong gian phòng này. Trọng lực ở đây thường xuyên được đẩy lên gấp 3 lần bình thường, còn nhiệt độ phòng thì duy trì ổn định ở mức\dots\ 45°C.''

Nodoka há hốc mồm: ``G-Gấp 3 lần trọng lực Trái Đất\dots!? Nhiệt độ 45°C!?''

Mặt đồng hồ trên tay Kuon lại lóe sáng xanh nhè nhẹ. Giọng Game cất lên chậm rãi, thả thêm một cú ``dội bom'' tinh thần vô cùng thản nhiên: ``Chưa hết đâu, Harusaki-san. Để rèn luyện phản xạ tuyệt đối cho các tài năng, dưới sàn nhà có tích hợp hệ thống cảm biến âm thanh siêu nhạy. Chỉ cần ai đó nhảy sai dù chỉ nửa nhịp của bài *Shiny Smily Story*, hệ thống phòng thủ tự động sẽ lập tức phóng ra sóng âm cao tần đánh văng người đó vào tường.''

Kuon khoanh tay trước ngực, gật gù xác nhận bằng giọng điệu không thể 'thật' hơn: ``Lần trước có một thành viên quên bài nhảy, bị sóng âm đánh văng làm thủng luôn mảng tường gạch chịu lực phía đông kìa. Anh vừa phải bỏ tiền túi gọi thợ đến vá lại tuần trước xong đấy.''

Mồ hôi hột túa ra như mưa trên trán Nodoka, hai chân cô bắt đầu run bẩy rẩy, mặt cắt không còn một giọt máu: ``T-Tại sao một phòng tập nhảy idol mà nghe cứ như phòng luyện tập siêu cấp của Bảy Viên Ngọc Rồng vậy ạ!? Mọi người tập nhảy hay là đang đi đánh trận thế!?''

A-chan bước tới vỗ vai Nodoka, ánh mắt tràn ngập vẻ ``thương hại'': ``Đó là cái giá của sự chuyên nghiệp em ạ. Một nhân viên quản lý như chúng ta cũng phải có thể lực phi thường để sẵn sàng xông vào môi trường trọng lực gấp 3 lần đó, vừa né sóng âm vừa đưa nước và khăn lau cho họ.''

Kuon chêm vào một câu chốt hạ: ``Nên là, nếu em không thể chịu đựng được áp lực trọng lực hoặc phản xạ không đủ nhanh để né sóng âm, em sẽ bị nghiền nát hoặc dính tường ngay khi vừa bước qua cánh cửa này đấy.''

Nodoka đứng chết trân tại chỗ. Cả không gian rơi vào im lặng vài giây.

Thế nhưng, đúng như kịch bản ``bất bình thường'' của cô tân binh này, thay vì sợ hãi xin lui, Nodoka đột ngột gập người một góc 90 độ, hai tay siết chặt thành nắm đấm, ngọn lửa quyết tâm trong mắt cô bùng cháy dữ dội hơn bao giờ hết: ``Em hiểu rồi ạ! Để có thể tiếp nước kịp thời và bảo vệ các tiền bối trong môi trường trọng lực khắc nghiệt đó, bắt đầu từ ngày mai em sẽ đeo tạ chân 10kg mỗi bên đi làm và đăng ký một khóa học võ thuật tự vệ để luyện phản xạ né sóng âm ạ!''

Nói rồi, cô nàng hừng hực khí thế lấy sổ tay ra, ghi chép lia lịa kế hoạch ``rèn luyện siêu nhân'' của mình.

Ba ``kẻ chủ mưu'' lập tức lùi lại vài bước, chụm đầu lại với nhau thì thầm trong sự bàng hoàng và bất lực toàn tập.

A-chan che miệng, thì thào với giọng run run: ``Kuon-san\dots\ chị bắt đầu thấy sợ cái sự quyết tâm của con bé rồi đấy\dots\ Nó định đeo tạ 20kg đi làm công sở thật kìa!''

Mặt đồng hồ của Kuon rung nhẹ, giọng Game phát ra thì thầm đầy lo ngại: ``Tôi e là nếu cứ tiếp tục cái kịch bản Cá tháng Tư này, đến cuối ngày chúng ta sẽ vô tình đào tạo ra một `siêu chiến binh công sở' có thể đấm nát cái trụ sở này mất\dots''

Kuon đỡ trán, thở dài một hơi đầy cay đắng: ``Cái trò đùa này đi quá xa rồi\dots\ Bây giờ mà bảo em ấy là `tất cả chỉ là cú lừa Cá tháng Tư', chắc em ấy đấm bọn mình bằng đôi tay đeo tạ thật quá\dots''

\ct

Chuyến tham quan tiếp tục diễn ra qua hàng loạt khu vực khác trong trụ sở, và kịch bản ``hù dọa'' của ba người họ cứ thế tiếp diễn với cấp độ ngày một leo thang.

Ở phòng máy chủ, A-chan úp úp mở mở về việc ``những oán hồn của các buổi livestream bị gỡ bỏ vẫn ẩn nấp trong các đường dây cáp'', Game tận tụy chêm vào bảng số liệu điện áp bất thường để tăng phần ma mị, thì Nodoka lại nghiêm túc rút sổ ghi chép kế hoạch mua muối mặn và bùa hộ mệnh về rải quanh các tủ server. Đến khu vực kho đạo cụ, Kuon cảnh báo rằng ``các bộ trang phục hóa trang có thể tự đứng dậy di chuyển vào giữa đêm'', thì Nodoka lại mắt sáng rỡ, thề sẽ học thêm khóa huấn luyện may vá và phong ấn đồ vật để ``quy phục'' chúng. Cứ mỗi lần ba người họ cố tung chiêu trò nắn gân, kết quả nhận lại luôn là một sự phản tác dụng hoàn toàn: Nodoka không những không chùn bước hay sợ hãi, mà còn đón nhận mọi thứ với một tinh thần sẵn sàng sinh tồn và hy sinh đến mức đáng kinh ngạc.

Cuối cùng, cả ba đưa cô tân binh trở lại gian văn phòng chính ban đầu. Bầu không khí trong phòng chùng xuống, trở nên vô cùng yên lặng.

A-chan khoanh tay trước ngực, hít một hơi sâu rồi quay sang nhìn đàn em: ``Được rồi, vậy là chuyến tham quan một vòng trụ sở cũng đã kết thúc. Em cảm thấy thế nào, Nodoka-san?''

Kuon tựa lưng vào mép bàn làm việc, ánh mắt điềm tĩnh quan sát cô gái trẻ: ``Có điều gì khiến em còn thắc mắc hay vẫn còn thấy sốc trước những gì vừa được chứng kiến không?''

Nodoka đứng đối diện hai vị tiền bối, cúi gằm mặt xuống. Cô ngập ngừng một hồi lâu, hai bàn tay siết chặt lấy tà áo công sở, cất giọng đầy trăn trở: ``\dots Thành thật mà nói, em thực sự cảm thấy rất lo lắng ạ\dots\ Bởi vì ở đây có quá nhiều thứ em chưa từng được biết tới, và có những điều quái gở đến mức em vẫn không tài nào hiểu nổi\dots''

Kuon khẽ mỉm cười dịu dàng, lên tiếng an ủi: ``Lần đầu đi làm ai mà chẳng bỡ ngỡ hả Nodoka-chan. Anh đây hồi mới chân ướt chân ráo bước vào cái văn phòng này cũng ngơ ngơ ngáo ngáo chẳng biết chỗ này vận hành kiểu gì. Thế mà rồi qua thời gian, cuối cùng đâu lại vào đấy hết ấy mà.''

Mặt đồng hồ trên tay Kuon lóe lên ánh sáng xanh nhè nhẹ. Giọng Game cất lên trầm ngâm, mang theo một lời đe dọa mang tính sinh tồn cuối cùng: ``Vấn đề cốt lõi là liệu cô có thể chịu đựng được những áp lực dị biệt và điên rồ sắp tới này hay không thôi. Nếu cô cảm thấy thể chất và tinh thần của mình không đủ kham nổi những hoạt động ngoài tầm kiểm soát của công ty này, thì rút lui về bờ an toàn ngay lúc này vẫn chưa muộn đâu, Harusaki-san.''

A-chan gật đầu, ánh mắt trở nên vô cùng nghiêm trọng: ``Đúng vậy. Bây giờ, em đang đứng trước một lựa chọn không thể đảo ngược lại. Nếu em muốn nộp đơn xin rút lui để tìm một công việc văn phòng bình thường khác, thì giờ chính là thời điểm thích hợp nhất--''

``\textbf{Nhưng mà!}''

Nodoka đột ngột ngẩng phắt đầu lên, cất tiếng hô to cắt ngang lời của các bậc tiền bối.

Ánh mắt cô tân binh không hề có một chút dao động hay sợ hãi. Trái lại, đôi mắt ấy đang rực cháy lên một niềm tin mãnh liệt. Cô bước tiến lên một bước, dõng dạc bày tỏ tâm tư từ tận đáy lòng: ``Đồng thời, tất cả những điều kỳ lạ ở đây đối với em đều là những trải nghiệm hoàn toàn mới mẻ và vô cùng thú vị! Tuy có phần nguy hiểm và khắc nghiệt, nhưng em thực sự cảm thấy \hll\ là một nơi làm việc siêu vui vẻ\dots\ Một nơi tràn đầy nhiệt huyết mà em có thể thấy bản thân mình thực sự trở thành một phần của nó!''

Lời nói tràn đầy sức nặng và sự chân thành của Nodoka vang vọng khắp căn phòng, khiến cả A-chan lẫn Kuon hoàn toàn sững sờ.

*“Con bé này… rốt cuộc là làm bằng thứ gì vậy chứ?”* - A-chan nhìn Nodoka mà trong lòng không khỏi chấn động. Dù vừa trải qua hàng loạt những lời hù dọa quái gở về nổ tung, vũ khí hạng nặng và cả trò ăn thịt đồng nghiệp, cô gái trẻ này không những không bỏ chạy mà còn nhìn thấy được ``nhiệt huyết'' bên trong đó.

*“Ánh mắt đó… hoàn toàn không phải là sự bốc đồng của một kẻ ngốc.”* - Kuon khẽ nhếch môi, sự nể phục dâng lên trong mắt cậu. Đứng trước vô vàn áp lực tưởng chừng như điên rồ nhất, thứ mà cô gái trước mặt lựa chọn không phải là từ bỏ, mà là dấn thân.

Nodoka siết chặt hai nắm đấm trước ngực, gập người một góc 90 độ chuẩn mực, hô lớn bằng tất cả sự quyết tâm: ``Vì vậy, dù có phải đối mặt với nổ tung hay trọng lực khắc nghiệt đi chăng nữa, em cũng sẽ cố gắng hết sức mình để giúp đỡ tất cả mọi người ở công ty này ạ!''

Không gian tĩnh lặng kéo dài vài giây. Nhìn cô tân binh tràn đầy sinh lực và lòng trung thành tuyệt đối đang cúi đầu trước mặt mình, A-chan và Kuon quay sang nhìn nhau. Cả hai không thể giữ nổi nét mặt nghiêm nghị được nữa, cùng bật cười nhẹ và nở một nụ cười mãn nguyện toàn tập.

A-chan mỉm cười ấm áp, lên tiếng: ``\dots Có vẻ như em thực sự rất phù hợp với vị trí này đấy, Nodoka-san.''

Kuon thở dài một hơi nhẹ nhõm, gãi gãi đầu: ``Và có lẽ\dots\ chúng ta cũng nên hạ màn kịch này xuống thôi nhỉ?''

Nodoka ngẩng đầu lên, ngơ ngác chớp mắt: ``Hể? Anh chị nói gì cơ ạ? Em không hiểu--''

*BỤP!*

Một tiếng nổ nhỏ vang lên giòn giã kèm theo cơn mưa tua rua giấy màu rực rỡ bắn tung tóe khắp không trung.

``Á!?'' Nodoka giật bắn mình, hai tay ôm lấy đầu hô lên.

Từ phía cánh cửa văn phòng, Tokino Sora bước vào với chiếc vỏ pháo bông nhỏ vừa được nổ trên tay. Cô nàng idol nở một nụ cười tươi thắm như hoa xuân, cất giọng dịu dàng: ``Mọi người vất vả nhiều rồi!''

Cùng lúc đó, Kuon tiến lại gần chiếc bảng trắng lớn ở góc phòng, thong thả lật mặt bảng lại. Trên mặt bảng lúc này hiện ra một bức vẽ bằng bút marker đen vô cùng ngộ nghĩnh mô tả lại hình ảnh A-chan đeo kính cận nghiêm túc trong trang phục văn phòng, xung quanh được trang trí bằng những sợi dây kim tuyến màu sắc rực rỡ từ từ thả xuống.

Nodoka ngơ ngác nhìn cô gái tóc nâu vừa bước vào, rồi lại nhìn bức vẽ trên bảng: ``K-Khoan đã\dots\ Chị là\dots!''

Sora, A-chan, Kuon và cả chiếc đồng hồ của Game cùng đồng thanh hô lớn: ``CHÚC MỪNG NGÀY CÁ THÁNG TƯ!''

Đôi mắt Nodoka sáng rực lên khi nhận ra người đứng trước mặt, nhưng rồi trí óc cô lập tức rơi vào trạng thái đình trệ: ``T-Tokino Sora-san đó sao!?\dots\ Mà khoan đã, ể!? C-Cá tháng Tư ư!?''

Sora nghiêng đầu cười tươi: ``Đúng vậy ạ! Hôm nay là ngày mùng 1 tháng 4 mà!''

Mặt đồng hồ của Kuon lóe sáng, giọng Game chêm vào một câu tỉnh rụi: ``Tôi nghĩ bởi vì cô quá háo hức và tập trung cho ngày đầu đi làm, nên đã hoàn toàn quên mất hôm nay là ngày gì rồi đúng không, Harusaki-san?''

Nodoka đứng chết trân, đầu óc quay mòng mòng như chong chóng. Cô lắp bắp, gương mặt ngập tràn sự hoang mang tột độ: ``À\dots\ À! Vậy\dots\ vậy tất cả những gì mà hai anh chị vừa dặn dò và chỉ bảo cho em từ sáng tới giờ\dots\ *bất chợt khựng lại, dường như vẫn đang bị ngợp*\dots\ Ủa? Thế cuối cùng cái nào là thật, và cái nào là trò đùa Cá tháng Tư vậy ạ!?''

Game phát ra một tiếng cười nhè nhẹ từ loa đồng hồ: ``Nói chung là cô sẽ sớm tự mình kiểm chứng thôi, Harusaki-san.''

Sora bước lại gần Nodoka, khẽ chắp hai tay trước ngực cười khúc khích: ``Cho em xin lỗi chị nhé, chắc lúc này chị đang bối rối lắm đúng không? Nhưng mà A-chan nói đúng đấy, chị thực sự là một người vô cùng chân thành và nghiêm túc.''

Kuon khoanh tay, mỉm cười gật gù đánh giá: ``Thú thật là tinh thần trách nhiệm và lòng quyết tâm của em làm bọn anh phải nể phục đấy. Không phải ai nghe xong mấy cái kịch bản điên rồ đó mà vẫn đòi đăng ký tập võ với đeo tạ đi làm như em đâu.''

Game từ trong chiếc đồng hồ cũng hùa theo: ``Chính xác. Sự kiên định của cô làm cái kịch bản dọa người của hai vị thư ký đây hoàn toàn phá sản.''

A-chan tiến lại gần, ánh mắt tràn đầy sự tin tưởng nhìn đàn em: ``Những đợt trước khi em hỗ trợ anh chị ở các dự án hợp tác khác, chị đã nhìn nhận em là một người rất chăm chỉ và chu đáo rồi.''

Kuon tiếp lời: ``Nhưng cũng chính vì bản tính quá thẳng thắn và thật thà của em, nên bọn anh nghĩ rằng trong ngày đầu tiên bước vào môi trường mới, em sẽ khá là lúng túng và căng thẳng.''

Sora mỉm cười tinh nghịch: ``Chính vì thế, nhân dịp đặc biệt này em cũng muốn phá vỡ một chút sự ngại ngùng ban đầu giữa mọi người, chị biết đấy ạ\~{}''

A-chan gãi má, hơi ngượng ngùng thú nhận: ``Bọn chị mới bày ra một trò chơi khăm nhỏ cho em nhằm nắn gân tinh thần một chút, và cũng là một cách để chúng ta gắn kết hơn trong những giây phút đầu tiên. Nhưng mà\dots''

Kuon thở dài, cúi đầu nhè nhẹ tỏ vẻ hối lỗi: ``\dots Dường như bọn anh có lẽ đã hơi làm quá lố rồi. Cho bọn anh xin lỗi em nhé, Nodoka-chan.''

Game nhóc nhách cất giọng trêu chọc từ chiếc đồng hồ: ``Cũng may là cô ấy đã dũng cảm nhận lời làm việc với chúng ta, chứ không thì chiều nay hai người có khi lại bị mấy anh cảnh sát tuýt còi mời lên phường uống trà vì tội ngược đãi và bạo hành nhân viên mới rồi đấy.''

A-chan lườm nhẹ chiếc đồng hồ trên tay Kuon, hắng giọng bào chữa: ``Chị cũng không thể ngờ rằng Sora lại đề nghị một trò chơi khăm quái chiêu như thế, và cũng chính Kuon-san là người bày ra mấy cái kịch bản nổ tung với ăn thịt dị bợm đó đấy nhé!''

Sora bật cười lém lỉnh, chỉ tay về phía Kuon: ``Kuon-senpai cũng đã nhiệt tình góp công lớn cho em còn gì nữa ạ!''

Game lắc đầu ngao ngán: ``Đúng là một cô Thần tượng lắm mưu nhiều mẹo. Còn Kuon-user thì cũng lầy lội chẳng kém cạnh gì.''

Kuon nhếch môi cười thản nhiên, nhún vai: ``Tất cả cũng chỉ để giúp cho em ấy cảm thấy thoải mái và tự nhiên nhất thôi mà, đúng không?''

Game lập tức bóc phốt: ``Tôi thì lại thấy các cô cậu làm cho Harusaki-san sợ tới mức hồn xiêu phách lạc thì có.''

A-chan bật cười, lắc đầu trước sự tung hứng của mọi người. Cô hắng giọng một tiếng thật rõ ràng để lấy lại bầu không khí chính thức: ``Dù sao đi nữa\dots\ Nodoka-san này\dots''

A-chan bước tiến lên một bước, nở một nụ cười vô cùng rạng rỡ và ấm áp, mìn nhận đàn em rồi chủ động chìa tay phải ra trước mặt Nodoka: ``Thay mặt cho ban quản lý và toàn bộ tập thể tài năng, cho phép anh chị chính thức chào mừng em tới với đại gia đình \hll!''

Nodoka nhìn bàn tay đang chìa ra trước mặt mình, rồi nhìn những nụ cười ấm áp của A-chan, Kuon, Sora và chiếc đồng hồ đang lóe sáng thân thiện của Game. Sự ngơ ngác trên khuôn mặt cô tân binh nhanh chóng tan biến, thay vào đó là một nụ cười hạnh phúc và tự hào trọn vẹn.

Nodoka vươn tay ra, siết chặt lấy bàn tay của A-chan, cúi đầu đầy trân trọng: ``Dạ! Em rất hân hạnh và vinh dự khi được làm việc với tất cả mọi người ạ!''

Ánh nắng buổi sáng ngày mùng 1 tháng Tư buông xuống qua ô cửa kính văn phòng, bao phủ lấy căn phòng trong sắc vàng ấm áp. Trò chơi khăm Cá tháng Tư đầy hỗn loạn và buồn cười đã khép lại, mở ra một chương mới rực rỡ tại trụ sở \hll, nơi thành viên nhân viên ban quản lý mới nhất mang tên Harusaki Nodoka chính thức bắt đầu hành trình của mình.

%==========================================TRUYỆN PHỤ=========================================
\omk

Khi tiếng cười giòn giã của ngày lễ Cá tháng Tư lắng xuống, A-chan chợt nhớ ra điều gì đó. Cô mỉm cười nhẹ nhàng, đưa tay lấy từ dưới gầm bàn ra một hộp quà hình vuông được bọc giấy trắng tinh khôi, thắt dải nơ màu đỏ thắm vô cùng tỉ mỉ rồi bước tới trước mặt Nodoka: ``Nhân tiện thì, chị có món quà này tặng cho em đây.''

Nodoka trố mắt nhìn hộp quà sang trọng trên tay A-chan, hai mắt sáng rỡ như sao sa, cô há hốc miệng vì kinh ngạc: ``Oa\dots\ Chiếc hộp đẹp quá! Đây có phải là quà chào mừng dành cho em không ạ!?''

Kuon gật gù, dựa lưng vào mép bàn, thong thả cất giọng: ``Thì, A-san với anh cất công chuẩn bị cho nhân viên mới thì cũng phải chăm chút một tí chứ. Làm sao có thể tặng cái gì đó xuề xòa được, đúng không?''

Nodoka ôm lấy hộp quà bằng cả hai tay, ánh mắt đầy sự mong chờ nhìn xung quanh: ``Vậy\dots\ anh chị sẽ không phiền nếu em mở nó ngay bây giờ chứ ạ?''

Mặt đồng hồ trên tay Kuon lóe lên ánh sáng xanh nhẹ, giọng Game vang lên điềm nhiên: ``Cứ việc. Quà của cô mà, Harusaki-san.''

Nodoka nghe vậy thì vô cùng háo hức. Cô đặt hộp quà lên bàn, ngón tay nhẹ nhàng rút dải nơ đỏ ra. Tiếng *PỰC* nho nhỏ vang lên cùng với giọng ngân nga *“Háo hức quá đi\~{}\~{}”* tràn đầy mong đợi của cô nàng tân binh.

Thế nhưng, trong lúc Nodoka đang mê mải gỡ nắp hộp, Sora, A-chan và Kuon lại quay sang nhìn nhau bằng những ánh mắt đầy vẻ ái ngại và tội lỗi. Cả ba người họ cùng chiếc đồng hồ trên tay Kuon đều biết rất rõ thứ gì đang nằm chờ chực bên trong chiếc hộp ấy.

Vừa gỡ chiếc nắp ra, nụ cười rạng rỡ trên môi Nodoka bỗng chốc đông cứng lại. Ánh mắt cô đờ đẫn nhìn xoáy vào bên dalam: ``Ủa? Cái gì đây ạ?\dots\ Trong này\dots\ có một núi tài liệu này?''

Bên trong chiếc hộp vuông vức không phải là bánh kẹo, hoa hay gấu bông, mà là từng chồng hồ sơ, báo cáo và giấy tờ công việc được xếp chồng chất cao ngất ngưởng như một ngọn núi nhỏ.

A-chan ho nhẹ một tiếng, đẩy gọng kính cận rồi cất giọng nghiêm túc của một vị quản lý chuyên nghiệp: ``Có thể em đúng là người mới thật, nhưng em cũng đã chính thức đặt chân vào công việc từ ngày hôm nay rồi đấy.''

Kuon xua xua tay tỏ vẻ hối lỗi nhưng nụ cười trên môi thì hoàn toàn không giấu nổi sự tỉnh rụi: ``Và chính vì thế, xin lỗi nếu điều này hơi khó nghe, nhưng anh chị sẽ bắt đầu giao việc cho em ngay lập tức, Nodoka-chan.''

Giọng Game từ trong chiếc đồng hồ phát ra một tiếng cười khàn nhẹ: ``Coi như là\dots\ bài tập khởi động nhẹ nhàng đi.''

Sora đứng bên cạnh chắp hai tay trước ngực, mỉm cười cổ vũ bằng tông giọng vô cùng ngọt ngào: ``Cố lên nhé, Nodoka-chan!''

Nodoka vẫn trố mắt nhìn đống giấy tờ, mồ hôi hột túa ra chảy dài trên má, giọng cô run run đầy hoang mang: ``Ể? E-Em có phải làm tất cả mọi thứ trong này không ạ?''

Kuon gật đầu cái rụp, trả lời không một chút do dự: ``Làm tất đấy em.''

Nodoka chớp chớp mắt, cố bấu víu lấy một tia hy vọng mỏng manh cuối cùng: ``Đ-Đây có phải lại là\dots\ một trò đùa Cá tháng Tư nữa không ạ?''

Mặt đồng hồ của Kuon rung nhẹ, Game cất giọng trầm ngâm đầy thương tiếc: ``Dù tôi rất muốn nói 'không phải' để cô đỡ sốc, nhưng rất tiếc, cái đống hồ sơ đó là thật 100% đấy.''

Nodoka gục đầu xuống, hai vai xuôi lơ: ``P-Phải ha\dots'' Thực tế phũ phàng của công việc công sở cuối cùng cũng giội thẳng một gáo nước lạnh lên đầu cô tân binh.

Nhìn vẻ mặt ỉu bảo của đàn em, A-chan dịu dàng bước lại gần, vỗ nhẹ lên vai cô an ủi: ``Em cũng không cần phải làm hết tất cả mọi thứ một mình đâu, Nodoka-san. Nếu có điều gì thắc mắc hay cần hướng dẫn thì cứ trao đổi trực tiếp với chị, Kuon-san, hoặc bất kỳ ai trong văn phòng này nhé.''

Kuon mỉm cười gật đầu chêm vào: ``Đúng đấy, buổi đầu làm việc nên lượng công việc giao cho em cũng sẽ tương đối nhẹ nhàng thôi, cứ yên tâm đi.''

Lời dặn dò chân thành của hai vị tiền bối như tiếp thêm một nguồn năng lượng khổng lồ. Ngọn lửa nhiệt huyết trong mắt Nodoka ngay lập tức bùng cháy trở lại. Cô đứng thẳng lưng, siết chặt hai nắm đấm, hô lớn bằng tất cả sự quyết tâm: ``\dots Dạ vâng ạ! Là thành viên mới nhất trong nhóm, em, Harusaki Nodoka, sẽ cố gắng hết sức để hoàn thành tốt mọi nhiệm vụ được giao!''

Sora nhìn cô tân binh hừng hực khí thế mà không khỏi trầm ồ, hai mắt sáng lên: ``Oa, Nodoka-chan đúng là tràn đầy năng lượng thật đấy! Nhìn em ấy quyết tâm như vậy làm em cũng thấy có thêm động lực rồi này!''

Kuon bật cười, nhấp một ngụm trà: ``Trông khí thế thế này thì chẳng mấy mà em ấy gánh vớt được phân nửa khối lượng công việc cho A-san rồi đấy chứ.''

A-chan chắp hai tay trước ngực, ánh mắt tràn ngập vẻ cảm động như vừa tìm thấy đấng cứu thế: ``Chính xác luôn! Có người san sẻ bớt đống báo cáo này thực sự là món quà lớn nhất đời chị lúc này rồi!''

Chiếc đồng hồ trên tay Kuon phát ra ánh sáng xanh nhè nhẹ, Game buông một câu nhận xét vừa hóm hỉnh vừa thực tế: ``Tôi chỉ hy vọng ngọn lửa nhiệt huyết rừng rực đó của cô ấy không bị đống giấy tờ ngập đầu kia dập tắt sau ba ngày làm việc đầu tiên thôi\dots''

Cả văn phòng lại được một trận cười rộn rã. Giữa bầu không khí ấm áp và tràn ngập niềm vui của ngày đầu tháng Tư, Nodoka hăng hái ôm đống tài liệu về chỗ ngồi của mình, sẵn sàng cho những thử thách mới. Và như vậy, câu chuyện chào đón nhân viên mới đầy hỗn loạn nhưng cũng vô cùng đáng nhớ của đại gia đình \hll\ chính thức khép lại!

\end{document}